\subsection*{Aufgabensammlung}

Hier sind Übungsaufgaben zur Wahrscheinlichkeitsrechnung. Deine Aufgabe ist
es, für jede Situation schnell zu entscheiden, ob du ein Baumdiagramm nutzt, die
Bernoulli-Formel (Binomialverteilung) anwendest oder die hypergeometrische
Verteilung (Ziehen ohne Zurücklegen aus einer endlichen Menge) benötigst.
\begin{enumerate}
\item
  In einer Urne liegen 10 rote und 5 blaue Kugeln. Es werden 4 Kugeln mit einem
Griff gezogen. Wie groß ist die Wahrscheinlichkeit, dass genau 2 rote Kugeln
dabei sind?\TRAINER{ 32.97\,\%}
\item
  Ein Multiple-Choice-Test besteht aus 10 Fragen mit jeweils 4
Antwortmöglichkeiten (nur eine ist richtig). Jemand rät bei jeder Frage. Wie
groß ist die Wahrscheinlichkeit, mindestens 8 Fragen korrekt zu beantworten?\TRAINER{ 0.0416\,\%} 
\item
  Eine Münze wird zweimal geworfen. Zeichne die möglichen Ausgänge auf und
berechne die Wahrscheinlichkeit für das Ereignis „Zuerst Zahl, dann Wappen“.\TRAINER{ 25\,\%}
\item
  In einer Schulklasse sind 15 Mädchen und 10 Jungen. Für ein Projekt werden 5
Personen zufällig ausgewählt. Wie hoch ist die Wahrscheinlichkeit, dass genau
3 Mädchen ausgewählt werden?\TRAINER{ Lotto: 195/506 ca. 38.54\,\%}
\item
  Ein Basketballspieler trifft seinen Freiwurf mit einer Wahrscheinlichkeit von 80\,\%. Er führt 12 Würfe aus. Wie groß ist die Wahrscheinlichkeit, dass er genau 10-
mal trifft?\TRAINER{ Bernoulli: 28.35\,\%}
\item
  Ein Glücksrad hat ein rotes und ein grünes Feld. Die Wahrscheinlichkeit für Rot
beträgt 0.3. Das Rad wird zweimal gedreht. Wie groß ist die Wahrscheinlichkeit
für die Kombination (1. Rot; 2. Grün)?\TRAINER{ 21\,\%} 
\item
  Aus einem Kartenspiel mit 32 Karten werden 3 Karten gleichzeitig gezogen.
Wie groß ist die Wahrscheinlichkeit, dass alle drei Karten Asse sind?\TRAINER{ ca 0.0806\,\% }
\item
  Bei einer Produktion von Schrauben sind erfahrungsgemäß 5\,\% der Teile
defekt. Einer Packung werden 20 Schrauben entnommen. Wie hoch ist die
Wahrscheinlichkeit, dass höchstens eine Schraube defekt ist?\TRAINER{ Bernoulli: 73.58\,\%}
\item
  Eine Wetterstation sagt für Samstag Regen mit 60\,\% und für Sonntag mit 40\,\%
  voraus. Mit welcher Wahrscheinlichkeit regnet es an beiden Tagen?\TRAINER{ 24\,\%}
  \newpage
\item
  In einer Obstschale liegen 6 Äpfel und 4 Birnen. Jemand nimmt mit
geschlossenen Augen 3 Früchte heraus. Wie groß ist die Wahrscheinlichkeit,
dass genau eine Birne dabei ist?\TRAINER{ Lotto: 50\,\%}
\item
  Ein Würfel wird 15-mal geworfen. Wie groß ist die Wahrscheinlichkeit, dass
  genau dreimal die Zahl „6“ gewürfelt wird?\TRAINER{ Bernoulli 23.63\,\%}
\item
  In einem Beutel sind 3 weiße und 7 schwarze Kugeln. Es wird eine Kugel
gezogen, die Farbe notiert und die Kugel beiseitegelegt. Dann wird eine zweite
Kugel gezogen. Wie groß ist die Wahrscheinlichkeit, zwei verschiedene Farben
zu erhalten?\TRAINER{ Lotto oder Baumdiagramm: 46.67\,\%}
\item
  Ein Unternehmen hat 50 Mitarbeiter, davon sind 10 im Betriebsrat. Es werden
6 Mitarbeiter für ein Interview zufällig bestimmt. Wie groß ist die
Wahrscheinlichkeit, dass kein Betriebsratsmitglied ausgewählt wird?\TRAINER{ Lotto 24.15\,\%}
\item
  Die Wahrscheinlichkeit für eine Mädchengeburt liegt bei etwa 51\,\%. Eine
Familie bekommt 4 Kinder. Wie groß ist die Wahrscheinlichkeit, dass
mindestens 2 davon Mädchen sind?\TRAINER{ Bernoulli: 70.23\,\%}
\item
  Ein Schütze trifft sein Ziel mit 90\,\%. Er gibt zwei Schüsse ab. Wie groß ist die
Wahrscheinlichkeit, dass er nur beim ersten Schuss trifft?\TRAINER{ Baum: 9\,\%}
\item
  Aus einer Kiste mit 20 Weinflaschen, von denen 4 verkorkt sind, werden 3
Flaschen für eine Feier geöffnet. Wie groß ist die Wahrscheinlichkeit, dass
genau eine der Flaschen verkorkt ist?\TRAINER{ Lotto: 42.11\,\%}
\item
  Ein Steuerberater stellt fest, dass 10\,\% aller eingereichten Belege fehlerhaft
sind. Er prüft 50 Belege stichprobenartig. Wie groß ist die Wahrscheinlichkeit,
dass genau 5 Belege Fehler aufweisen?\TRAINER{ Bernoulli 18.49\,\%}
\item
  Ein Jäger schießt auf eine Ente. Die Trefferwahrscheinlichkeit beträgt 0.7. Er
schießt maximal zweimal. Wenn er beim ersten Mal trifft, hört er auf. Mit
welcher Wahrscheinlichkeit erlegt er die Ente erst mit dem zweiten Schuss?\TRAINER{ 21\,\%}

\end{enumerate}
\newpage
